% Shared project plan table. Included by every weekly report and by the thesis appendix.
% Needs: tabularx, booktabs (xcolor optional) -- inherited from the including document.

\section*{Master thesis objectives}

The overall goal of this MSc thesis is to develop and validate a computer-vision framework that
segments the coronary artery tree from coronary computed tomography angiography and
characterizes its topology, and to apply it to a large population cohort. The project
encompasses the following main objectives:

.

\begin{itemize}
  \item Describe the function of the heart with a special focus on the role of coronary arteries
  \item Describe how computed tomography is used to record and visualize the coronary arteries
  \item Describe the common topology of the coronary arteries and its anatomical variants.
  \item Review the state of the art in coronary artery segmentation and in topological
    description of the coronary arteries.
  \item Implement, validate and document a state-of-the-art segmentation framework for
    coronary arteries from computed tomography.
  \item Develop, implement, revise and validate methods for topological characterization of
    coronary artery trees.
  \item Extract topological features such as branching pattern, dominance, bifurcation
    angles and tortuosity.
  \item Apply the framework to a large cohort and produce population descriptions and
    statistics.
  \item Implement and test outlier detection and extreme-value detection of the topological
    features.
  \item Associate the topological features with patient outcome.
\end{itemize}






\subsection*{Project plan}

Please find the too long project plan on the next page :)) 

\begin{table}[htbp]
  \centering
  \begin{tabularx}{\textwidth}{@{}l X r@{}}
    \toprule
    \textbf{Week} & \textbf{Activity} & \textbf{Risk} \\
    \midrule
    1 & Project kick-off; ImageCAS-X and the existing code base explored; clinical
        background started & 1 \\
    2 & Clinical background continued: anatomy, dominance, variants, CAD and CCTA; topology
        research; literature review and state of the art (obj.\ 1--4); dataset chapter for
        ImageCAS-X & 1 \\
    \midrule
    \multicolumn{3}{@{}l}{\textit{Track A --- ImageCAS-X, in parallel with Track B}} \\
    \cmidrule(r){1-3}
    3--4 & Understand in depth the deep-learning segmentation method chosen.
           Topology feature extraction on ImageCAS-X: branching, bifurcation angles,
           tortuosity and dominance. Feature reliability against the ground truth for the best
           segmentation method, on the 160 test cases (obj.\ 5--7) & 3  \\
    5    & Topology-aware fine-tuning targeted at the failure modes found in week 3 from the delivered weights, using a
           method from the literature review chosen to address them, if the gap is large
           enough to justify it; re-evaluated on the same 160 test cases & 3 \\
    5--6 & Preparing the move to the CGPS cohort: ground-truth-free inference path, and
           per-scan connectivity quality control & 2 \\
    \midrule
    \multicolumn{3}{@{}l}{\textit{Track B --- CGPS enabling, in parallel with Track A}} \\
    \cmidrule(r){1-3}
    2--3  & CGPS data access and exploration: repeat-scan intervals, outcome statistics, major adverse cardiac events and more & 5 \\
    4  & Training in the in-house labelling software; labelling protocol agreed & 2 \\
    5--6 & Annotation of a small subset of the CGPS subset & 3 \\
    \midrule
    \multicolumn{3}{@{}l}{\textit{Merged --- CGPS deployment and analysis (Multiple ideas to follow)}} \\
    \cmidrule(r){1-3}
    7--10   & Deploy the segmentation and topology pipeline on CGPS: inference, feature
             extraction and per-scan quality control & 4 \\
    7-10 & Accuracy validation against the labelled subset; ImageCAS-X to CGPS domain
             shift & 3 \\
    7--10 & Method for anatomical correspondence between repeat scans of the same
             participant & 4 \\
    7--10 & Repeat-scan reproducibility study: per-feature measurement error and minimum
             detectable change (obj.\ 6) & 3 \\
    7--10 & Population description of the cohort and normative reference ranges for the
             topological features (obj.\ 8) & 3 \\
    11--14 & Outlier and extreme-value detection on the topological features (obj.\ 9) & 3 \\
    11--14 & Association of baseline geometry with plaque progression, at locations free of
             plaque at baseline (obj.\ 10) & 5 \\
    11--14 & Risk-prediction model for major adverse cardiac events (obj.\ 10): a baseline model
             on traditional and clinical risk factors from CGPS, and a second model adding
             topological features --- both continuous measurements and the population-outlier
             flags from the extreme-value analysis above (obj.\ 9, e.g.\ top-5\% tortuosity) ---
             to test whether topology improves prediction beyond the baseline & 5 \\
    15--16 & Thesis writing, discussion and limitations, abstract and so on & 2 \\
    22 December & Report submission & 2 \\
    5 January    & Oral presentation & 1 \\
    \bottomrule
  \end{tabularx}
  \caption{Project plan for the development of my thesis :)) }  
  \label{tab:project_plan}
\end{table}

